% !TEX root = ../aa_SoM_Chapters.tex

%% HARRIS: Chapter 10

% ö = \%c3\%b6
% ä = \%C3\%A4

\xdef\sectiontitle{\IIsectiontitle} 

%%%%%%%%%%%%%%%
\begin{frame}[label=2_8_a]%[label=2_7_TOC1]
\frametitle{\texorpdfstring{\culine[\sectiontitleulinecolor]{\sectiontitle}}{\sectiontitle} }

%\vspace{0.5cm}
%\begin{itemize}[leftmargin=0.5cm,itemsep=3mm]
%    \item[2.1] \hyperlink{2_1_a}{\IIaSubsectiontitle}
%    \item[2.2] \hyperlink{2_2_a}{\IIbSubsectiontitle}	
%    \item[2.3] \hyperlink{2_3_a}{\IIcSubsectiontitle}	
%    \item[2.4] \hyperlink{2_4_a}{\IIdSubsectiontitle}
%    \item[2.5] \hyperlink{2_5_a}{\IIeSubsectiontitle}
%    \item[2.6] \hyperlink{2_6_a}{\IIfSubsectiontitle}
%    \item[2.7] \hyperlink{2_7_a}{\CDAlert[blue]{\IIgSubsectiontitle}}
%    \item[2.8] \hyperlink{2_8_a}{\IIhSubsectiontitle}
%    \item[2.9] \hyperlink{2_9_a}{\IIiSubsectiontitle}
%\end{itemize}

\vspace{0.2cm}
\begin{itemize}[leftmargin=0.5cm,itemsep=3mm]
    \item[2.1] \hyperlink{2_1_a}{\IIaaSubsectiontitle}
    \item[2.2] \hyperlink{2_2_a}{\IIaSubsectiontitle}
    \item[2.3] \hyperlink{2_3_a}{\IIbSubsectiontitle}	
    \item[2.4] \hyperlink{2_4_a}{\IIcSubsectiontitle}	
    \item[2.5] \hyperlink{2_5_a}{\IIdSubsectiontitle}
    \item[2.6] \hyperlink{2_6_a}{\IIeSubsectiontitle}
    \item[2.7] \hyperlink{2_7_a}{\IIfSubsectiontitle}
    \item[2.8] \hyperlink{2_8_a}{\CDAlert[blue]{\IIgSubsectiontitle}}
    \item[2.9] \hyperlink{2_9_a}{\IIhSubsectiontitle}
    \item[2.10] \hyperlink{2_10_a}{\IIiSubsectiontitle}
\end{itemize}


\endofslide \end{frame}
%%%%%%%%%%%%%%%

\xdef\sectiontitle{\IIgSubsectiontitle} 

%%%%%%%%%%%%%%%
\begin{frame}
\frametitle{\texorpdfstring{\culine[\sectiontitleulinecolor]{\sectiontitle}}{\sectiontitle} }

\framesubtitle{Binary compound semiconductors}
\seminormal

%Compound semiconductors
\begin{itemize}[itemsep=1.7mm]
	\item Elemental semiconductors are mostly the elements in group IV: \appear{Si, Ge}
	
\item \CDAlert[fuchsia]{Compound semiconductors}, i.e. mostly mixtures of III-V\appear{, II-VI} \appear{and IV-IV elements}\appear{, are also widely used} \appear{(also Si$_{1-x}$Ge$_{x}$)}
\appear{\xdef\loc{south east}%
\begin{tikzpicture}[remember picture,overlay] % anchor=south
    \node (A) [yshift=-0.5*\figYshift,xshift=\figXshift, anchor=\loc] at (current page.\loc) {\includegraphics[page=1,height=5.8cm]{Figures_booklet/SOM_semiconductor_table_figure.png}};%scale=\figscale. % 8cm max height
\end{tikzpicture}}
\item Varying of the concentration of the components is used to tailor the properties of the semiconduc-\\ 
\end{itemize}
% 6.5 cm = midway, 10 cm = 3/4 
%\vspace{2.5mm}
\begin{minipage}{5.1cm}
	\begin{itemize}[itemsep=1.7mm]
	\item[] tors (most notably the band gap\appear{, and the curvatures of the bands)}

\item Figure shows the \Alert[red]{band gap \appear{vs. lattice constant}} \appear{for some common binary compound semiconductors\footnote{www.tf.uni-kiel.de} }

\item This tunability is useful when forming functional circuit components and devices!
\end{itemize}

\end{minipage}

% great WWW: https://www.tf.uni-kiel.de/matwis/amat/semi_en/kap_5/backbone/r5_1_4.html
%\endofslide 
%\endofsection
\end{frame}
%%%%%%%%%%%%%%%

%%%%%%%%%%%%%%%
\begin{frame}
\frametitle{\texorpdfstring{\culine[\sectiontitleulinecolor]{\sectiontitle}}{\sectiontitle} }
\framesubtitle{Diode junction}
%The Diode
 
 \seminormal
\begin{itemize}
	\item Perhaps the most common semiconductor device is a \CDAlert[fuchsia]{diode} \appear{(or a diode junction)}\appear{, allowing electrical \Alert[fuchsia]{current to flow only along one direction} }
	
	\item A diode is assembled using blocks of p-type \appear{and n-type doped semiconductors}\appear{, 'glued' on atomic level together with some appropriate impurity element}\appear{, in overall forming a so-called \CDAlert[fuchsia]{p–n junction}}	
	
	\item If we apply a potential difference with the p-type at higher potential\appear{, known as \Alert[red]{forward-biased condition}}\appear{, electrons flow to the left and  holes to the right, both toward the junction }
\end{itemize}
	
% 6.5 cm = midway, 10 cm = 3/4 
\vspace{2.5mm}
\begin{minipage}{6.5cm}
	\begin{itemize}	
	\item When they meet, they recombine \implies{ \Alert[red]{current flows continuously} } 

\end{itemize}

\end{minipage}

\xdef\loc{south east}
\begin{tikzpicture}[remember picture,overlay] % anchor=south
    \node (A) [yshift=-0.25*\figYshift,xshift=\figXshift, anchor=\loc] at (current page.\loc) {\includegraphics[page=1,width=8cm]{Figures_booklet/SOM_10_Bonding_figures/SOM_Bonding_Figure_41.png}}; %scale=\figscale. % 8cm max height
\end{tikzpicture}

\endofslide 
%\endofsection
\end{frame}
%%%%%%%%%%%%%%%

%%%%%%%%%%%%%%%
\begin{frame}
\frametitle{\texorpdfstring{\culine[\sectiontitleulinecolor]{\sectiontitle}}{\sectiontitle} }
\framesubtitle{Diode junction}
%The Diode
 
 \seminormal
\begin{itemize}
	\item One of the most common semiconductor devices is \CDAlert[fuchsia]{diode} (or diode junction), allowing electrical \Alert[fuchsia]{current to flow only along one direction} 
	
	\item A diode is assembled using blocks of p-type and n-type doped semiconductors, 'glued' on atomic level together with some appropriate impurity element, in overall forming a so-called \CDAlert[fuchsia]{p–n junction}	
	 
	
	\item The other bias is to connect high potential to $n$-type side\appear{, i.e. \CDAlert[blue]{reverse-biased condition} }
	\implies{ Both electrons and holes move\\
		      away from the junction, and\\
		      \Alert[blue]{current stops almost instantly}  }
		      
	\item Note that this behavior is very\\ 
		different from e.g. resistors\appear{, that\\
		comply with the Ohm's law}
\end{itemize}
	
\xdef\loc{south east}
\begin{tikzpicture}[remember picture,overlay] % anchor=south
    \node (A) [yshift=-0.25*\figYshift,xshift=\figXshift, anchor=\loc] at (current page.\loc) {\includegraphics[page=1,width=8cm]{Figures_booklet/SOM_10_Bonding_figures/SOM_Bonding_Figure_41.png}}; %scale=\figscale. % 8cm max height
\end{tikzpicture}

\endofslide 
%\endofsection
\end{frame}
%%%%%%%%%%%%%%%

%%%%%%%%%%%%%%%
\begin{frame}
\frametitle{\texorpdfstring{\culine[\sectiontitleulinecolor]{\sectiontitle}}{\sectiontitle} }
\framesubtitle{p-n junction}
%p-n junction

\begin{itemize}[itemsep=1.5mm]
	\item Thus, we connect two doped semiconductors\appear{, let's say p-type Ge on the left and n-type Ge on the right} \appear{to make a \CDAlert[fuchsia]{p–n junction}}
	
	\item When they are separate\appear{, they have different Fermi levels $E_F$} 
	
	\item After connecting, electrons 'high' on the n-type side cross to the p-type
\end{itemize}

% 6.5 cm = midway, 10 cm = 3/4 
%\vspace{2.5mm}
\begin{minipage}{5.4cm}
	\begin{itemize}[itemsep=1.5mm]
	\item[] side, and holes from p-type cross to the n-side 
	
	\item After a short migration of carriers to either sides\appear{, the system finds a new macroscopic equilibrium state}
	
	\item In equilibrium, the Fermi levels will be equal on both sides \appear{forming an energy level scheme shown in Figure (a)}
	
%	\item In addition, there will be a small build up of negative and positive charges into either sides of the depletion zone
\end{itemize}

\end{minipage}

\appear{\xdef\loc{south east}
\begin{tikzpicture}[remember picture,overlay] % anchor=south
    \node (A) [yshift=-0.25*\figYshift,xshift=1*\figXshift, anchor=\loc] at (current page.\loc) {\includegraphics[page=1,height=5.8cm]{Figures_booklet/SOM_10_Bonding_figures/SOM_Bonding_Figure_42.png}}; %scale=\figscale. % 8cm max height
\end{tikzpicture}}

\endofslide 
%\endofsection
\end{frame}
%%%%%%%%%%%%%%%

%%%%%%%%%%%%%%%
\begin{frame}
\frametitle{\texorpdfstring{\culine[\sectiontitleulinecolor]{\sectiontitle}}{\sectiontitle} }
\framesubtitle{Reverse-biased p-n junction}

\begin{itemize}[itemsep=1.5mm]
	\item The energy level diagrams in Fig.~42 clarify the behaviour of the diode.
 \appear{In the unbiased case (a)}\appear{, in fact, only few electrons need to cross the junction until an equilibrium is achieved}
	
	\item The immediate region of the junction (the impregnated semiconductor 
\end{itemize}

\begin{minipage}{5.3cm}
	\begin{itemize}[itemsep=1.5mm]
	\item[] paint or glue), called \Alert[red]{depletion zone}, is free of charge carriers \appear{(it is too wide to tunnel through)}
	
	\item In the case of \Alert[blue]{reverse-biasing} (b, the difference of the energy bands is exaggerated)\appear{, minority carriers (the opposite sign carriers) are scarce}

\end{itemize}
\end{minipage}

\xdef\loc{south east}%
\begin{tikzpicture}[remember picture,overlay] % anchor=south
    \node (A) [yshift=-0.25*\figYshift,xshift=\figXshift, anchor=\loc] at (current page.\loc) {\includegraphics[page=1,height=5.8cm]{Figures_booklet/SOM_10_Bonding_figures/SOM_Bonding_Figure_42.png}};%scale=\figscale. % 8cm max height
\end{tikzpicture}
\endofslide 
%\endofsection
\end{frame}
%%%%%%%%%%%%%%%

%%%%%%%%%%%%%%%
\begin{frame}
\frametitle{\texorpdfstring{\culine[\sectiontitleulinecolor]{\sectiontitle}}{\sectiontitle} }
\framesubtitle{Forward-biased p-n junction}

\begin{itemize}[itemsep=1.5mm]
	\item The energy level diagrams in Fig.~42 clarify the behaviour of the diode. In the unbiased case (a), in fact, only few electrons need to cross the junction until an equilibrium is achieved
	
	\item The immediate region of the junction (the impregnated semiconductor %paint or glue), called depletion zone, is free of charge carriers. It is too wide to tunnel through. In the case of reverse-biasing (b), the difference of the energy bands is exaggarated, and minority carriers (the opposite sign carriers) are scarce
\end{itemize}

%\vspace{2.5mm}
\begin{minipage}{5.3cm}
	\begin{itemize}[itemsep=1.5mm]
	\item[] paint or glue), called \Alert[red]{depletion zone}, is free of charge carriers (it is too wide to tunnel through) 
	
	\item The \Alert[red]{forward-biasing} (c) shifts the $\mathrm{n}$-type conduction band upward %relative to that of the $\mathrm{p}$-type 
	\implies{ Electrons and holes near the\\ 
		      \CDAlert[red]{depletion zone} diffuse across it\appear{, \\%diffuse cross
		      recombine and induce a \CDAlert[red]{current} }}
\end{itemize}
\end{minipage}

\xdef\loc{south east}
\begin{tikzpicture}[remember picture,overlay] % anchor=south
    \node (A) [yshift=-0.25*\figYshift,xshift=\figXshift, anchor=\loc] at (current page.\loc) {\includegraphics[page=1,height=5.8cm]{Figures_booklet/SOM_10_Bonding_figures/SOM_Bonding_Figure_42.png}}; %scale=\figscale. % 8cm max height
\end{tikzpicture}

\endofslide 
%\endofsection
\end{frame}
%%%%%%%%%%%%%%%

%%%%%%%%%%%%%%%
\begin{frame}
\frametitle{\texorpdfstring{\culine[\sectiontitleulinecolor]{\sectiontitle}}{\sectiontitle} }
\framesubtitle{Currents inside the depletion zone}%of the p-n junction

\seminormal%
\begin{itemize}[itemsep=1.5mm]
	\item There are actually four different currents across the junction: \appear{two currents} \appear{(drift and diffusion)} \appear{for both \CDAlert[red]{majority carriers}} \appear{and two for the \CDAlert[blue]{minority carriers}} 
	\appear{\xdef\loc{south east}%
\begin{tikzpicture}[remember picture,overlay] % anchor=south
    \node (A) [yshift=-0.25*\figYshift,xshift=1*\figXshift, anchor=\loc] at (current page.\loc) {\includegraphics[page=2,width=5cm]{Figures_booklet/Tikz_SOM_Bonding_figures_pn_junction}};%
\end{tikzpicture}%
}

	\item Electrons are swept opposite to the field to the n-side, and holes are swept correspondingly to the p-side
	
	\item Diffusion of states then leads to \CDAlert[red]{recombination currents} for holes and electrons\appear{, denoted by $i_{\text {holes,rec}} \Equiv i_{p, r}$} \appear{and $i_{\text {electrons,rec}} \Equiv i_{n, r}$}

\item Electron–hole pairs are also generated near the\\
	junction due to thermal excitation%\appear{, and are swept through the junction due to the build-up E-field in the \CDAlert[fuchsia]{depletion region} }

 
\end{itemize}

\vspace{1.5mm}
\begin{minipage}{8.45cm}
	\begin{itemize}[itemsep=1.5mm]
\item The corresponding currents\appear{, called \Alert[blue]{(thermal)} \CDAlert[blue]{generation currents}, can be labelled by $i_{\text {holes,gen}} \Equiv i_{p,g} \quad$} \appear{and $\quad i_{\text {electrons,gen}} \Equiv i_{n,g}$}

\item At equilibrium\appear{, the generation and recombination currents are equal:} %
$
 \appear{| i_{p, g} |}  \equal \appear{| i_{p, r} |} \appear{~and~} \appear{| i_{n, g} |}  \equal  \appear{| i_{n, r} |}
$
\end{itemize}
\end{minipage}


\endofslide 
%\endofsection
\end{frame}
%%%%%%%%%%%%%%%

%% LEAVE DIFFUSION REGIONS AND THE DEPLETION REGION TO FOLLOW-UP COURSES??
%%%%%%%%%%%%%%%%
%\begin{frame}
%\frametitle{\texorpdfstring{\culine[\sectiontitleulinecolor]{\sectiontitle}}{\sectiontitle} }
%
%\begin{itemize}
%	\item Taking a closer look at the junction area, There are actually three zones:
%
%\item Depletion region in the middle, across the junction
%
%\item Diffusion regions on both sides of the depletion region, from there the electrons and holes diffuse to the junction area
%
%\item Homogeneous regions far away from Home the junction, which are actually not much affected by the junction
%\end{itemize}
%
%%\xdef\loc{south east}
%%\begin{tikzpicture}[remember picture,overlay] % anchor=south
%%    \node (A) [yshift=-0.25*\figYshift,xshift=1*\figXshift, anchor=\loc] at (current page.\loc) {\includegraphics[page=6,height=4cm]{Figures_booklet/Tikz_SOM_Bonding_figures}}; %scale=\figscale. % 8cm max height
%%\end{tikzpicture}
%
%\endofslide 
%%\endofsection
%\end{frame}
%%%%%%%%%%%%%%%


%%%%%%%%%%%%%%%
\begin{frame}
\frametitle{\texorpdfstring{\culine[\sectiontitleulinecolor]{\sectiontitle}}{\sectiontitle} }
\framesubtitle{Currents in p-n junction}

\seminormal%

\begin{itemize}[itemsep=1.5mm]
	\item Now we apply \CDAlert[fuchsia]{forward bias}\appear{, i.e. positive potential difference across the junction }
	\appear{\xdef\loc{south east}%
\begin{tikzpicture}[remember picture,overlay] % anchor=south
    \node (A) [yshift=-0.25*\figYshift,xshift=1*\figXshift, anchor=\loc] at (current page.\loc) {\includegraphics[page=3,width=5cm]{Figures_booklet/Tikz_SOM_Bonding_figures_pn_junction}};%
\end{tikzpicture}%
}
	\item A forward bias decreases the electric field in the depletion region\appear{, as well as it decreases the differences between the energy levels on the p- and n-sides by an amount of:}
\[
	\Delta E  \equal \appear{-e \cdot V}
\]
\appear{It becomes easier for electrons to diffuse to $p$-region} \appear{and for the holes to $n$-region.} \appear{This effect \CDAlert[fuchsia]{increases both recombination currents} by the Boltzmann factor:}
%$$
%e^{-\Delta E / k_{B} T}=e^{e \cdot V / k_{B} T}
%$$
%(We don't need Fermi-Dirac statistics, since most of the available states of diffusing electrons and holes are empty) 
\end{itemize}

% 6.5 cm = midway, 10 cm = 3/4 
\vspace{2.5mm}
\begin{minipage}{8.45cm}
	\begin{itemize}[itemsep=1.5mm]
\item[] \[
\appear{e^{-\Delta E / k_{B} T}}  \equal \appear{e^{e \cdot V / k_{B} T}}
\]
\appear{where use of Fermi–Dirac statistics is omitted}\appear{, since most of the available states of diffusing electrons and holes are empty}

\item A current from p- to n-side is generated
\end{itemize}
\end{minipage}

\endofslide 
%\endofsection
\end{frame}
%%%%%%%%%%%%%%%

%%%%%%%%%%%%%%%%
%\begin{frame}
%\frametitle{\texorpdfstring{\culine[\sectiontitleulinecolor]{\sectiontitle}}{\sectiontitle} }
%
%\begin{itemize}
%	\item Now we apply forward bias, i.e. positive potential difference across the junction. A forward bias decreases the electric field in the junction region, as well as it decreases the differences between the energy levels on the $p$ - and $n$-sides by an amount of:
%$$
%\Delta E=-e \cdot V
%$$
%It becomes easier for electrons to diffuse to $p$-region and for the holes to $n$-region. This effect increases both recombination currents by the Boltzmann factor:
%$$
%e^{-\Delta E / k_{B} T}=e^{e \cdot V / k_{B} T}
%$$
%(We don't need Fermi-Dirac statistics, since most of the available states of diffusing electrons and holes are empty).
%\end{itemize}
%
%%\xdef\loc{south east}
%%\begin{tikzpicture}[remember picture,overlay] % anchor=south
%%    \node (A) [yshift=-0.25*\figYshift,xshift=\figXshift, anchor=\loc] at (current page.\loc) {\includegraphics[page=1,width=7cm]{Figures_booklet/SOM_10_Bonding_figures/SOM_Bonding_Figure_42.png}}; %scale=\figscale. % 8cm max height
%%\end{tikzpicture}
%
%\endofslide 
%%\endofsection
%\end{frame}
%%%%%%%%%%%%%%%%


%%%%%%%%%%%%%%%
\begin{frame}
\frametitle{\texorpdfstring{\culine[\sectiontitleulinecolor]{\sectiontitle}}{\sectiontitle} }
\framesubtitle{Currents in p-n junction}

\begin{itemize}
	\item Now we can 'calculate' the total current for e.g. the holes:
\[
i_{p, t o t}   \equal \appear{i_{p, r}}  \minus \appear{| i_{p, g} | }  \equal  \appear{| i_{p, g} | }\appear{e^{e \cdot V / k_{B} T}}  \minus \appear{ | i_{p, g} |}  \equal  \appear{| i_{p, g} | }\appear{(e^{e \cdot V / k_{B} T}-1)}
\]
\appear{The same calculation can be performed also for electrons}\appear{, thus we get the total current:}
\[
 I  \equal \appear{i_{p, t o t}}  \plus \appear{i_{n, t o t}} \equal \appear{I_{S}(T) ( e^{e \cdot V / k_{B} T}-1 )}
 \]
\appear{where $I_{S}$ is called the \Alert[red]{(reverse bias)} \CDAlert[red]{saturation current}}\appear{, and equation is known as the \CDAlert[fuchsia]{Shockley diode equation} }

\item Discussion can be repeated for reverse bias as well\appear{, and it turns out that the same equation is valid for both positive and negative values }

\item The obtained I–V curve, representing the behaviour of a typical \CDAlert[fuchsia]{diode rectifier}\appear{, showing $I=I(V)$}\appear{, \Alert[fuchsia]{differs markedly from Ohm's law}}
\end{itemize}

\endofslide 
%\endofsection
\end{frame}
%%%%%%%%%%%%%%%

%%%%%%%%%%%%%%%
\begin{frame}
\frametitle{\texorpdfstring{\culine[\sectiontitleulinecolor]{\sectiontitle}}{\sectiontitle} }
\framesubtitle{I–V curve of the p-n junction}

\begin{itemize}
	\item By plotting the current $I(V)$\appear{, given by the \CDAlert[fuchsia]{Shockley diode equation}}
\[
I(V)  \equal  \appear{I_{S}(T)} \appear{( e^{e \cdot V / k_{B} T}-1 )}
\]
\appear{as a function of external voltage $V$}\appear{, we see that the behavior \CDAlert[fuchsia]{differs markedly from Ohm's law}}
\appear{\xdef\loc{south east}%
\begin{tikzpicture}[remember picture,overlay] % anchor=south
    \node (A) [yshift=-0.25*\figYshift,xshift=\figXshift, anchor=\loc] at (current page.\loc) {\includegraphics[page=1,width=7cm]{Figures_booklet/SOM_Tipler_Solid_state_physics_figures/SOM_Tipler_SSP_figure_46}};%scale=\figscale. % 8cm max height
\end{tikzpicture}}
\end{itemize}

\vspace{2.5mm}
\begin{minipage}{6.7cm}
	\begin{itemize}%[itemsep=1.5mm]
\item This is a characteristic behaviour of a diode\appear{, and is difficult to achieve using conventional electronic components} \appear{(resistors}\appear{, capacitors} \appear{and inductors)}

\item At very large reverse-biased voltage values \appear{(\CDAlert[blue]{breakdown voltage} = threshold)}\appear{, the pn-junction starts to pass current in the opposite direction}
\end{itemize}
\end{minipage}

\endofslide 
%\endofsection
\end{frame}
%%%%%%%%%%%%%%%

%%%%%%%%%%%%%%%
\begin{frame}
\frametitle{\texorpdfstring{\culine[\sectiontitleulinecolor]{\sectiontitle}}{\sectiontitle} }
\framesubtitle{Light-emitting diode (LED)}

\seminormal
\begin{itemize}[itemsep=1.7mm]
	\item \CDAlert[red]{Light-emitting diode (LED) is a p-n junction that emits light }
	
	\item When the junction is forward biased\appear{, many holes are pushed from their $\mathrm{p}$-region to the junction region}\appear{, and many electrons are pushed from $n$-region to the junction region }
	\implies{ Many electrons recombine with holes in the junction region }

\item During recombination\appear{, the electron can emit a photon with energy approximately equal to the band gap }

\end{itemize}

\vspace{1.7mm}
\begin{minipage}{8.7cm}
	\begin{itemize}[itemsep=1.7mm]
	\item This energy \appear{(i.e. colour of emitted light)} \appear{can be varied by using materials with different band gaps (see earlier slide) }

\item LEDs are widely used for lighting \appear{and for digital displays in electronic equipments} \appear{etc.}

\item  A diode of band gap $1$\;eV would need a forward bias of $1$\;V before turning on
\end{itemize}
\end{minipage}

\xdef\loc{south east}%
\begin{tikzpicture}[remember picture,overlay] % anchor=south
    \node (A) [yshift=-0.25*\figYshift,xshift=1*\figXshift, anchor=\loc] at (current page.\loc) {\includegraphics[page=5,width=5cm]{Figures_booklet/Tikz_SOM_Bonding_figures_pn_junction}};%
\end{tikzpicture}%

\endofslide 
%\endofsection
\end{frame}
%%%%%%%%%%%%%%%

%%%%%%%%%%%%%%%
\begin{frame}
\frametitle{\texorpdfstring{\culine[\sectiontitleulinecolor]{\sectiontitle}}{\sectiontitle} }
\framesubtitle{Photodiode}

\seminormal
%A Photodiode
\begin{itemize}[itemsep=2.2mm]
	\item Consider an unbiased diode with band gap of $1$\,eV\appear{, illuminated  with photons whose energies slightly exceed $1$\;eV} 
	\item \Alert[red]{Junction absorbs photons in all regions lifting electrons into the conduction band} (and holes into the valence band) %A conduction electron in the p-side (where it is a minor carrier) is eager to slide down to the lower conduction band in the entype. Similarly holes created in either the depletion zone or $\mathrm{n}$-type valence zone would happily float up to the p-type valence. Thus, as long as the photons arrive from an external source, the component acts like a battery,
\end{itemize}

\vspace{2.2mm}
\begin{minipage}{7.8cm}
	\begin{itemize}[itemsep=2.2mm]
	%\item[] a hole in the valence band. 
	
	\item Conduction \CDAlert[red]{electron} in the p-side \appear{(where it is a minor carrier)} \appear{is eager to \CDAlert[red]{slide down} to the lower conduction band in the n-side} 
	
	\item Similarly \CDAlert[blue]{holes} created in either the depletion zone or n-side valence zone\appear{, would happily \CDAlert[blue]{float up} to the p-side}
	
	\item Thus, as long as the photons arrive from an external source\appear{, the component acts like a battery} \appear{(\Alert[fuchsia]{creates a voltage that generates a current})}
	\end{itemize}
\end{minipage}

\xdef\loc{south east}
\begin{tikzpicture}[remember picture,overlay] % anchor=south
    \node (A) [yshift=-0.15*\figYshift,xshift=\figXshift, anchor=\loc] at (current page.\loc) {\includegraphics[page=1,width=5.7cm]{Figures_booklet/SOM_10_Bonding_figures/SOM_Bonding_Figure_43.png}}; %scale=\figscale. % 8cm max height
\end{tikzpicture}

\endofslide 
%\endofsection
\end{frame}
%%%%%%%%%%%%%%%

%%%%%%%%%%%%%%%
\begin{frame}
\frametitle{\texorpdfstring{\culine[\sectiontitleulinecolor]{\sectiontitle}}{\sectiontitle} }
\framesubtitle{Photovoltaic cell}

\seminormal
\begin{itemize}[itemsep=1.8mm]
	\item If the band gap is $1$\;eV\appear{, then the photon energy must be at least the band gap energy:}
\[
\lambda \equal \frac{h c}{E_{gap}}  \equal \frac{1240\;eV \cdot nm}{1\; eV} \equal \appear{1240\;nm}
\]

\item Connecting this device to an external circuit\appear{, it becomes a source of emf and power }
\implies{ Such a device is often called a \CDAlert[fuchsia]{solar cell or photovoltaic cell} } 

\item Sunlight is not necessarily required\appear{, since photon energies greater than the band gap will suffice}

\item Here we treated the essence of the photovoltaic cell\appear{, which turns light into energy}\appear{, providing the tools for \CDAlert[fuchsia]{solar energy conversion}} %We may also add layers of intrinsic semiconductor between the n-type and $p$-type regions, widening the depletion zone where the electron-hole pairs are sorted quickly and contributing to improved efficiengy in high-frequency applications

\item But the same physics applies also to other light-sensitive devices \appear{(i.e. \Alert[red]{photodiodes}):} \appear{charge-coupled device (CCD) image detectors}\appear{, complementary metal–oxide–semiconductor (\CDAlert[fuchsia]{CMOS}) detectors}\appear{, digital cameras} \appear{and camcorders}
\end{itemize}

\endofslide 
%\endofsection
\end{frame}
%%%%%%%%%%%%%%%

%%%%%%%%%%%%%%%
\begin{frame}
\frametitle{\texorpdfstring{\culine[\sectiontitleulinecolor]{\sectiontitle}}{\sectiontitle} }
\framesubtitle{Diode laser}
%The Diode Laser

\seminormal
\begin{itemize}[itemsep=1.7mm]
	\item As in the LED, the energy given up in electron–hole recombination \appear{is the source of light in a \CDAlert[fuchsia]{diode laser} }
	
	\item Large concentrations of electrons and holes at the pn-junction \appear{provide the condition of \CDAlert[fuchsia]{population inversion}}\appear{, required for lasing action}

	\item A resonant \CDAlert[fuchsia]{cavity} is also needed to amplify the desired wavelength, but %due to the high refractive indices of many semiconductors (specifically for Si, $n \sim 3.7$), semiconductor–air interfaces often act as good enough cavity mirrors
\end{itemize}

%\vspace{1.7mm}
\begin{minipage}{8.2cm}
\begin{itemize}[itemsep=1.7mm]
\item[] due to the high refractive indices of many semiconductors \appear{(e.g. for Si, $n\,{\sim}\,3.7$)}\appear{, semiconductor–air interfaces often act as good enough cavity mirrors}

\item \CDAlert[red]{Indirect band gap Si} is not the optimal choice for diode lasers\appear{, since it easily gives up its recombination energy in forms other than light}

\item \CDAlert[blue]{GaAs} is a III-V semiconductor\appear{, and a good material due it its \CDAlert[blue]{direct band gap} }

	\end{itemize}
\end{minipage}

\xdef\loc{south east}
\begin{tikzpicture}[remember picture,overlay] % anchor=south
    \node (A) [yshift=-0.25*\figYshift,xshift=\figXshift, anchor=\loc] at (current page.\loc) {\includegraphics[page=1,width=5.3cm]{Figures_booklet/SOM_Tipler_Solid_state_physics_figures/SOM_Tipler_SSP_figure_52.png}}; %scale=\figscale. % 8cm max height
\end{tikzpicture}

\endofslide 
%\endofsection
\end{frame}
%%%%%%%%%%%%%%%

%%%%%%%%%%%%%%%
\begin{frame}
\frametitle{\texorpdfstring{\culine[\sectiontitleulinecolor]{\sectiontitle}}{\sectiontitle} }
\framesubtitle{Transistor}
% The Transistor

\begin{itemize}[itemsep=1.7mm]
	\item Transistors in their simplest configuration consist of three extrinsic semiconductor blocks\appear{, either in \CDAlert[fuchsia]{npn}} \appear{or pnp configuration (see Figure)}
	\appear{\xdef\loc{south east}%
\begin{tikzpicture}[remember picture,overlay] % anchor=south
    \node (A) [yshift=-0.25*\figYshift,xshift=1*\figXshift, anchor=\loc] at (current page.\loc) {\includegraphics[page=1,width=13cm]{Figures_booklet/SOM_Tipler_Solid_state_physics_figures/SOM_Tipler_SSP_figure_53.png}};%scale=\figscale. % 8cm max height
\end{tikzpicture}}
%\item Here npn is discussed

\item The \Alert[red]{p-type in the middle is the \CDAlert[red]{base}}\appear{, and the \Alert[blue]{n-type blocks are the \CDAlert[blue]{emitter}} \appear{and the \CDAlert[blue]{collector}}}

\item Emitter and collector are of different size and dopant levels\appear{, tailored for their role}
\end{itemize}

\endofslide 
%\endofsection
\end{frame}
%%%%%%%%%%%%%%%

%%%%%%%%%%%%%%%
\begin{frame}
\frametitle{\texorpdfstring{\culine[\sectiontitleulinecolor]{\sectiontitle}}{\sectiontitle} }
\framesubtitle{npn-junction transistor}

\begin{itemize}[itemsep=1.7mm]
	\item \CDAlert[blue]{Forward-biased emitter} diode allows conduction electrons to flow from emitter into the base

\item Thin and \CDAlert[red]{lightly-doped base} is p-type
	\implies{ electrons are minority carriers, but do not likely recombine with holes
 }
\end{itemize}
\begin{minipage}{5cm}
\begin{itemize}[itemsep=1.7mm]
\item Collector is an n-type semiconductor\appear{, and the \Alert[blue]{collector–base diode is reverse biased} }

\item Small fraction of electrons recombine with holes in the collector\appear{, while majority 'slide down' to the collector's conduction band}\appear{, creating a \Alert[fuchsia]{current which can be controlled with the base}}
	\end{itemize}
\end{minipage}

\xdef\loc{south east}
\begin{tikzpicture}[remember picture,overlay] % anchor=south
    \node (A) [yshift=-0.25*\figYshift,xshift=\figXshift, anchor=\loc] at (current page.\loc) {\includegraphics[page=1,width=8.5cm]{Figures_booklet/SOM_10_Bonding_figures/SOM_Bonding_Figure_44.png}}; %scale=\figscale. % 8cm max height
\end{tikzpicture}

\endofslide 
%\endofsection
\end{frame}
%%%%%%%%%%%%%%%

%%%%%%%%%%%%%%%
\begin{frame}
\frametitle{\texorpdfstring{\culine[\sectiontitleulinecolor]{\sectiontitle}}{\sectiontitle} }
\framesubtitle{Transistor amplifier}

\begin{itemize}
	\item The fraction of the electrons that flows out the collector \appear{(eventually returning to the emitter through the output circuit)} \appear{is \CDAlert[red]{proportional} to} %and typically 100 times greater than that which flows out the base (returning through the input circuit).
\end{itemize}

\begin{minipage}{6.6cm}
\begin{itemize}%[itemsep=1.7mm]
\item[]  and \Alert[red]{typically 100 times greater than the signal flowing into the base} 

\item In Fig.~45,%
\appear{\xdef\loc{south east}%
\begin{tikzpicture}[remember picture,overlay] % anchor=south
    \node (A) [yshift=-0.25*\figYshift,xshift=\figXshift, anchor=\loc] at (current page.\loc) {\includegraphics[page=1,width=6.8cm]{Figures_booklet/SOM_10_Bonding_figures/SOM_Bonding_Figure_45.png}}; %scale=\figscale. % 8cm max height
\end{tikzpicture}%
}%
the input introduces a minute voltage variations in the emitter-base bias\appear{, with proportional variations in the rate at which electrons flow into the base from the emitter.} \appear{A tiny fraction of this current leaves the base and circulates as a varying current in the input (emitter-base) circuit}
	\end{itemize}
\end{minipage}

\endofslide 
%\endofsection
\end{frame}
%%%%%%%%%%%%%%%

%%%%%%%%%%%%%%%
\begin{frame}
\frametitle{\texorpdfstring{\culine[\sectiontitleulinecolor]{\sectiontitle}}{\sectiontitle} }
\framesubtitle{Field-effect transistor (FET)}

\semismall
\begin{itemize}[itemsep=1.4mm]
	\item Very common switching devices in \CDAlert[fuchsia]{integrated circuits} are \Alert[red]{field-effect transistors (FETs)} 
	
	\item A FET can be formed into a p-type silicon slab\appear{, by introducing two n-type regions to the top}\appear{, called '\CDAlert[blue]{source}' and '\CDAlert[tuniblue]{drain}'}
	
	\item  A third electrode\appear{, called '\CDAlert[red]{gate}'} \appear{is separated from the slab}\appear{, source and drain by an insulating layer of silicon dioxide ($\mathrm{SiO}_{2}$)} 	
		
\end{itemize}

\vspace{1.4mm}
\begin{minipage}{7.3cm}
\begin{itemize}[itemsep=1.4mm]
	\item With uncharged gate\appear{, negligible current passes through the drain and source}\appear{, because one of the pn-junctions is reverse biased} 
	
	\item But for a positively-charged gate\appear{, already small charge} \appear{(because gate is minuscule)} \appear{provides a substantial electric field}\appear{, which enables \CDAlert[red]{current to flow through the npn-junction}  }
	
	\item FETs are tiny \appear{(${\sim}\,1\;\mu m$)} \appear{integrated transistors}
	
	\item \Alert[red]{MOSFET}{https://en.wikipedia.org/wiki/MOSFET\#Importance} = Metal–oxide–semiconductor FET
	\end{itemize}
\end{minipage}

\xdef\loc{south east}
\begin{tikzpicture}[remember picture,overlay] % anchor=south
    \node (A) [yshift=-0.25*\figYshift,xshift=\figXshift, anchor=\loc] at (current page.\loc) {\includegraphics[page=1,width=6.5cm]{Figures_booklet/SOM_10_Bonding_figures/SOM_Bonding_Figure_46.png}}; %scale=\figscale. % 8cm max height
\end{tikzpicture}

%\endofslide 
\endofsection
\end{frame}
%%%%%%%%%%%%%%%

%%%%%%%%%%%%%%%%
%\begin{frame}
%\frametitle{\texorpdfstring{\culine[\sectiontitleulinecolor]{\sectiontitle}}{\sectiontitle} }
%
%\begin{itemize}
%	\item A single-electron transistor (SET) is a sensitive electronic device based on the Coulomb blockade effect
%
%\item In this device the electrons flow through a tunnel junction between source/drain to a quantum dot (conductive island)
%\end{itemize}
%
%%\xdef\loc{south east}
%%\begin{tikzpicture}[remember picture,overlay] % anchor=south
%%    \node (A) [yshift=-0.25*\figYshift,xshift=\figXshift, anchor=\loc] at (current page.\loc) {\includegraphics[page=1,width=7cm]{Figures_booklet/SOM_10_Bonding_figures/SOM_Bonding_Figure_42.png}}; %scale=\figscale. % 8cm max height
%%\end{tikzpicture}
%
%\endofslide 
%%\endofsection
%\end{frame}
%%%%%%%%%%%%%%%%
